\section{Introduction}\label{sec:introduction}

Synthetic financial time series are increasingly studied for applications including portfolio optimisation and volatility forecasting \citep{hounwanou2025evaluating}, as well as the evaluation of trading strategies under alternative simulated paths \citep{zhang2026sfag}.
The reliability of these applications depends on whether generated series preserve the statistical structures that govern financial risk and market dynamics.
Financial returns exhibit persistent empirical regularities, commonly referred to as stylized facts, including heavy-tailed marginal distributions, volatility clustering, and aggregational Gaussianity \citep{cont2001empirical}.
In particular, the autocorrelation of absolute returns can remain positive over long lags, providing a quantitative signature of volatility persistence \citep{ding1993long}.
A generator that reproduces one of these properties while distorting another may produce plausible-looking data but still yield misleading conclusions in downstream financial analyses.

Evaluating this fidelity remains difficult.
Existing approaches commonly assess whether synthetic observations can be distinguished from real ones \citep{yoon2019timegan}, whether predictive models trained on synthetic data transfer to real data \citep{esteban2017realvalued}, or whether the real and synthetic samples are close under a global distributional distance \citep{gretton2012kernel}.
These approaches provide useful overall assessments of realism or utility, but they generally do not identify which financially relevant structures are preserved or distorted.
Finance-specific evaluations instead commonly report whether generated series recover selected stylized facts, such as heavy tails, volatility clustering, leverage effects, or gain/loss asymmetry \citep{takahashi2019modeling, tanaka2025cofindiff}.
Although such checks are financially meaningful, demonstrating that a property is present does not by itself establish whether its magnitude, persistence, or conditional structure matches that of the real data.

This limitation is particularly important because recent generators increasingly incorporate financial properties directly into their training objectives.
\citet{zhang2026sfag}, for example, formulate differentiable losses that encourage generated series to reproduce selected stylized facts.
When a property has been explicitly optimized during training, merely confirming its presence provides limited independent evidence of fidelity.
Evaluation must instead determine how closely the magnitude, persistence, and conditional structure of the generated property match those of real data.
More generally, recent work has documented specific failure modes of generative-data metrics and advocated targeted sanity checks for determining whether a metric detects known deficiencies \citep{raisa2025position}.

To address this need, we propose a diagnostic framework for evaluating synthetic financial return panels.
Rather than reducing fidelity to a single global measure, the framework evaluates four financially relevant properties separately: the tail-weighted marginal distribution, long-lag volatility clustering, aggregational Gaussianity, and regime-conditional tail shape.
Each synthetic--real discrepancy is calibrated against the disagreement between matched real-data subsamples, providing a common reference for interpreting otherwise heterogeneous metrics.
The resulting profile identifies which structures a generator preserves and where it departs from real financial data.

% RUN4-VERIFIED-INTRODUCTION: only AIL M4 includes 0.5.
{\color{red}{\color{black}
We apply the framework to one-minute returns from 600 U.S.-listed
equities and compare geometric Brownian motion, multi-scale stochastic
volatility, AIL, and Schr\"odinger--Bass Bridge for Time Series (SBBTS).}
The results distinguish their statistical profiles: AIL's marginal
distribution, volatility-clustering, and aggregational-Gaussianity
intervals fall below the empirical reference of 0.5, while its
regime-conditional-tail interval includes it.
The comparison uses 200 outer ticker-corpus replicates, with a separate
20-seed analysis of numerical stability on the fixed observed panels.
An exploratory 19-specification analysis over five paired seeds further
examines sensitivity to the selected M1 and M4 parameters; the observed
orderings are retained across these alternatives.
{\color{black}SBBTS improves on both controls for the latter two properties but remains
close to GBM on the first two.}
{\color{black}The main contribution is therefore a financially interpretable evaluation
protocol that localizes discrepancies relative to the variability observed
within real data.
}}
